\documentclass[10pt]{article}

\usepackage[a4paper,margin=0.62in]{geometry}
\usepackage[T1]{fontenc}
\usepackage[utf8]{inputenc}
\usepackage{enumitem}
\usepackage{tabularx}
\usepackage{xcolor}
\usepackage{microtype}
\usepackage[
  unicode=true,
  colorlinks=true,
  linkcolor=black,
  citecolor=black
]{hyperref}

\definecolor{ink}{HTML}{1F2937}
\definecolor{muted}{HTML}{5F6B7A}
\definecolor{rule}{HTML}{D9DEE8}
\definecolor{accent}{HTML}{5B3F8C}
\hypersetup{urlcolor=accent}

\setlength{\parindent}{0pt}
\setlength{\parskip}{0pt}
\setlength{\emergencystretch}{2em}
\renewcommand{\familydefault}{\sfdefault}

\newcommand{\rssection}[1]{%
  \vspace{0.85em}
  {\large\bfseries\color{accent} #1}\par
  \vspace{0.2em}
  {\color{rule}\hrule height 0.6pt}
  \vspace{0.45em}
}

\newcommand{\bodypar}[1]{%
  {\normalsize #1\par}
  \vspace{0.58em}
}

\begin{document}
\pagestyle{empty}
\color{ink}

\begin{tabularx}{\linewidth}{@{}Xr@{}}
  {\LARGE\bfseries Research Statement} &
  {\small\color{muted}Updated July 2026} \\[0.48em]
  {\Large\bfseries Junhuai Xu} &
  {} \\
  {\normalsize Ph.D. Candidate, Department of Physics, Tsinghua University} &
  {} \\
  {\normalsize \textbf{Email}: \href{mailto:xjh22@mails.tsinghua.edu.cn}{xjh22@mails.tsinghua.edu.cn}} &
  {} \\
  {\normalsize \textbf{Ph.D. Advisor}: \href{https://inspirehep.net/authors/1062622}{Zhigang Xiao}} &
  {} \\
  {\normalsize \href{https://orcid.org/0000-0001-7529-8786}{ORCID} \quad \href{https://inspirehep.net/authors/1866896}{INSPIRE-HEP} \quad \href{https://scholar.google.com/citations?user=XK-LVeIAAAAJ\&hl=en}{Google Scholar}} &
  {}
\end{tabularx}

\vspace{0.65em}
{\color{accent}\hrule height 1.0pt}
\vspace{0.8em}

\rssection{Research Theme}
\bodypar{My research lies at the intersection of experimental nuclear and particle physics, detector instrumentation, and scientific computing. A central theme across my work is the extraction of hidden physical information from complex experimental observables. In heavy-ion collisions, this hidden information may appear as high-momentum nucleons generated by short-range correlations or as freeze-out source distributions encoded in two-particle correlation functions. In neutrino experiments, it may appear as event topology, direction, energy, and oscillation information embedded in sparse detector responses.}

\bodypar{These problems share a common structure: experiments measure detector-filtered, statistically fluctuating observables rather than the physical quantities of interest directly. My work develops physics-informed reconstruction and inference methods that connect detector-level measurements to robust physics conclusions. This includes detector calibration, response simulation, event reconstruction, inverse-problem formulation, Bayesian and likelihood-based inference, machine-learning reconstruction, and uncertainty propagation.}

\rssection{Short-Range Correlations and Hard-$\gamma$ Probes}
\bodypar{One major direction of my Ph.D. research is the study of nucleon short-range correlations (SRCs) in nuclei. Traditional mean-field descriptions treat nucleons as moving predominantly below the Fermi momentum, but strong short-distance nucleon--nucleon interactions generate correlated nucleon pairs and high-momentum components beyond the mean-field picture. My work uses hard bremsstrahlung $\gamma$ rays emitted in low-energy heavy-ion collisions as a clean electromagnetic probe of these high-momentum nucleons.}

\bodypar{In the CSHINE hard-$\gamma$ program, high-energy photons from ${}^{124}$Sn+${}^{124}$Sn collisions at 25 MeV/u are measured with the CSHINE-Gamma CsI(Tl) hodoscope. The key physical idea is that SRC-induced high-momentum nucleons increase the relative energy of early-stage neutron-proton collisions, hardening the emitted $np\rightarrow np\gamma$ spectrum. By comparing precision $\gamma$-ray measurements with detector-filtered IBUU-MDI transport calculations, I extracted the SRC-induced high-momentum-tail fraction in ${}^{124}$Sn as \(R_{\mathrm{HMT}}=(20\pm3)\%\). Under the adopted HMT parametrization, this corresponds to substantial SRC-related high-momentum components in both neutron and proton distributions.}

\bodypar{A significant part of this work was experimental and methodological rather than only phenomenological. I contributed to CSHINE-Gamma calibration and response validation, including high-energy CsI(Tl) response studies with quasi-monochromatic photon beams at SLEGS/SSRF. I developed analysis workflows for rare bremsstrahlung $\gamma$-ray signals, including event reconstruction, background control, detector-response corrections, systematic uncertainty evaluation, and Richardson--Lucy spectrum unfolding as an independent validation of the main forward-folding analysis. This experience shaped my broader view that reliable physics extraction depends on treating detector response and uncertainty propagation as central parts of the measurement.}

\rssection{Femtoscopic Imaging and Inverse Problems}
\bodypar{A second research direction is correlation-function imaging in relativistic heavy-ion collisions. Two-particle correlation functions contain femtoscopic information about the freeze-out source, but the source distribution is not directly observable. Conventional analyses often assume a Gaussian source and extract a small number of radius parameters. My work instead treats femtoscopy as an inverse problem: reconstruct the real-space source distribution from the measured correlation function without imposing a predefined source shape.}

\bodypar{I developed a Richardson--Lucy-based imaging framework that combines Hanbury Brown--Twiss femtoscopy, the Koonin--Pratt equation, and final-state-interaction models. Applied to STAR proton-proton and antiproton-antiproton correlations in Au+Au collisions at \(\sqrt{s_{NN}}=200\) GeV, this method reconstructed non-Gaussian freeze-out source functions for matter and antimatter. The proton and antiproton source functions were found to be consistent within uncertainties, providing coordinate-space evidence for matter--antimatter symmetry at freeze-out. The reconstructed sources also reveal non-Gaussian structures that would be hidden by a purely Gaussian ansatz.}

\bodypar{This work demonstrates how a method originally developed for optical deblurring can be adapted into a physics-constrained inference framework. The response kernel is not generic; it is determined by two-particle wave functions and final-state interactions. The reconstruction must preserve positivity, remain stable against statistical fluctuations, and be validated through pseudo-data and resampling tests. These requirements connect nuclear femtoscopy with a broader class of inverse problems in experimental physics.}

\newpage

\rssection{Neutrino Simulation and Machine-Learning Reconstruction}
\bodypar{In parallel with my heavy-ion research, I have worked on atmospheric-neutrino simulation and reconstruction for the TRIDENT deep-sea neutrino telescope program. This direction extends my experience from dense detector arrays and $\gamma$-ray spectroscopy to sparse, large-volume neutrino detection. I built an atmospheric-neutrino simulation chain combining HONDA fluxes, FLUKA interactions, and Geant4 detector transport for compact underwater detector configurations.}

\bodypar{On the reconstruction side, I developed graph-neural-network methods to infer event class, direction, and energy from detector responses. This is a natural setting for machine learning because the detector observations are sparse, geometrically structured, and noisy. At the same time, the analysis must remain grounded in physics: reconstruction performance must be evaluated with controlled simulation samples, detector geometry, optical propagation, and oscillation-parameter sensitivity studies. My work therefore emphasizes simulation-to-inference workflows in which ML models are connected to physics observables and uncertainty evaluation rather than used as isolated prediction tools.}

\rssection{Methodological Foundation}
\bodypar{Across these projects, my methodological interest is in physics-informed reconstruction. I am particularly interested in problems where the measurement can be expressed as a response kernel acting on a hidden physical distribution, with statistical fluctuations and systematic uncertainties layered on top. In this setting, reconstruction is not simply a numerical inversion; it requires physical modeling of the response, stable iterative or probabilistic inference, validation with closure tests, and uncertainty propagation to the final observable.}

\bodypar{The methods I have used include Richardson--Lucy deconvolution and imaging, likelihood-based inference, pseudo-experiment studies, response-matrix construction, Geant4 detector simulation, resampling-based uncertainty evaluation, graph neural networks, and high-performance computing workflows. I have also explored sequential Bayesian methods such as Unscented Kalman Filtering and particle filtering in partially observed dynamical systems. These experiences reinforce a common methodological direction: interpretable inference tools should make their assumptions explicit and should be tested against controlled simulations and detector-level effects.}

\rssection{Future Directions}
\bodypar{My long-term goal is to develop interpretable AI-for-science frameworks for experimental nuclear and particle physics. I am especially interested in reconstruction methods that combine machine learning with physical constraints, response modeling, and statistically meaningful uncertainty estimates. In the near term, I plan to extend this agenda along three connected directions.}

\begin{itemize}[leftmargin=1.35em,itemsep=0.25em,topsep=0.1em]
  \item \textbf{Physics-informed inverse problems:} developing robust reconstruction and unfolding methods for spectra, source functions, and detector-filtered observables, with emphasis on validation, stability, and uncertainty propagation.
  \item \textbf{ML-based event reconstruction:} applying graph-based and geometry-aware learning methods to sparse detector responses in neutrino and nuclear experiments, while connecting model outputs to downstream physics inference.
  \item \textbf{Unified detector-to-physics workflows:} building analysis pipelines that integrate calibration, detector simulation, reconstruction, systematic studies, and data-model comparison in reproducible computational frameworks.
\end{itemize}

\vspace{0.45em}
\bodypar{The broader motivation is to make stronger and more transparent physics statements from increasingly complex experimental data. Future nuclear and particle experiments will rely not only on more capable detectors, but also on reconstruction and inference frameworks that respect detector physics, encode prior knowledge carefully, and quantify uncertainty honestly. I aim to contribute as a researcher who can work across these boundaries: from detector response and experimental analysis to inverse problems, statistical inference, and machine-learning reconstruction.}

\end{document}
